\documentclass{article}
\usepackage{fontspec}
\usepackage{amssymb}
\usepackage{graphicx}
\usepackage{sectsty}
\usepackage{enumitem}
\usepackage{tikz}  % for drawing boxes
\usepackage[margin=1in]{geometry}  % for managing margins

\newcommand{\bigbox}{\raisebox{\depth}{\scalebox{1.5}{$\square$}}}

\setmainfont{Roboto}

\sectionfont{\large}  % make section titles normal size

\begin{document}

% Large box for iteration number
\noindent\begin{tikzpicture}
\draw (0,0) rectangle (10,2);
\node at (5,1) {\huge Iteration Number: };
\end{tikzpicture}

% Smaller boxes for experiment name and description
\noindent\begin{tikzpicture}
\draw (0,0) rectangle (7,1);
\node at (2,0.5) {\textbf{Experiment Name: }};
\end{tikzpicture}
\hspace{1cm}  % horizontal space between the boxes
\begin{tikzpicture}
\draw (0,0) rectangle (7,1);
\node at (2,0.5) {\textbf{Description: }};
\end{tikzpicture}

\section{Computational Requirements for Bayesian Optimization}
\begin{itemize}
\item[\bigbox] {Generate samples using Sobol sampling}
\item[\bigbox] {Generate the layout for \textbf{Opentrons II robot} transfer, generate a tick box list for the manual mixes, and the automated csv file for the \textbf{Triple Quadrupole mass spectrometer}}

This will generate 5 files: A master table with the run information, a Opentrons II robot script, an QqQ csv table for mass spec run, a diagram with the microplate and the corresponding condition, and a table with the microplate layout and the colors.

\item[\bigbox] {Generate tables for the microtubes mixes}
\\\\
\textbf{Pause!}

\item[\bigbox] {Dataset prepared and preprocessed for the optimization task}
\item[\bigbox] {Give new recommendations using the acquisition function}

This will give files:
\item[\bigbox] {Come back to the second step in a new sheet, all the boxes should be ticked}
\end{itemize}

\section{Opentrons II Robot Preparation}
\begin{itemize}
\item[\bigbox] {Test the programmed protocol with water and dummy plates}
\\\\
\textbf{On the preparation day!}

\item[\bigbox] {Check the stock solutions that are needed for the experiment. Prepare fresh ones when required. Keep it fridged. Filter-sterilise when it's possible, specially with concentrated solutions.}
\item[\bigbox] {Prepare the necessary mixes and keep in the fridge}
\item[\bigbox] {Calibration of the robot (very important)}
\item[\bigbox] {Loading of necessary labware (pipette tips, plates, etc.)}
\item[\bigbox] {Programming the experimental protocol (it's given before)}
\item[\bigbox] {Loading of samples and reagents}
\item[\bigbox] {Run the programmed protocol}
\\\\
\textbf{After samples are collected} 

\item[\bigbox] {Run the programmed protocol to make the QC pool}

\end{itemize}

\section{Plate Reader Preparation}
\begin{itemize}
\item[\bigbox] {Prepare agar plate with the bacteria (one day before)}
\item[\bigbox] {Prepare complete media for pre-cultures}
\item[\bigbox] {Prepare precultures (5ml complete media)}
\item[\bigbox] {Calibration of the plate reader \textbf{(Currently not done by user)}}
\item[\bigbox] {Setting up the main script with the correct parameters (wavelength, time, temperature, etc.)}
\item[\bigbox] {Book the slot for the plater reader in the building system}
\\\\
\textbf{On the day of the experiment!}

\item[\bigbox] {Setting up the software with the preheating programme (15 minutes before)}
\item[\bigbox] {Preparing and loading the sample plate (Come from \textbf{Opentrons II robot} }

\item[\bigbox] {Verifying the functioning and cleanliness of the reading chamber  \textbf{(Currently not done by user)}}
\item[\bigbox] {Running the plate}
\\\\
\textbf{On between!}
\item[\bigbox] {Prepare ice before taking off the plates}
\item[\bigbox] {Weight and label every tube that will be used}

\textbf{After 24 / 36 / 48 hrs}
\item[\bigbox] {Take off the plate from the reader}
\item[\bigbox] {Put the plate on ice, mix the wells with the tip, and move 500 ul to the tubes (that are on ice \textbf{Keep the plates for biofilm assay})}
\item[\bigbox]{Quenching with cold methanol/water, the ratio should be }
\item[\bigbox] {Centrifuge the tubes in a cold centrifuge}
\item[\bigbox] {Take 400ul of the tubes into new ones cold, remove the rest of supernatant, keep everything on ice \textbf{Weight the pellet quickly}}
\item[\bigbox] {Move the tubes to the -80 freezer}
\end{itemize}

\section{Triple Quadrupole Mass Spectrometer Preparation}
\begin{itemize}
\item[\bigbox] {Calibration of the mass spectrometer\textbf{(Currently not done by user)}}
\item[\bigbox] {Cleaning and verifying the ion source and detector \textbf{Has to be done once a while} }
\item[\bigbox] {Open csv for the sequence list, create new sequence from this file}
\item[\bigbox] {Check that there's enough mobile phase solvent for the run and miliQ water for the seal wash}
\item[\bigbox] {Check the status of the mass spec by looking at the parameters, and if they are in range}

\item[\bigbox] {Run a few blanks, plus system suitability sample, plus a few blanks}
\item[\bigbox] {Preparing and loading the samples from the experiment}
%\item[\bigbox] {Optimizing the mass spectrometer settings for the analytes of interest}
\item[\bigbox] {Perform the run}
\end{itemize}

% Dashed line to indicate a break
%\noindent\textbf{\rule{\textwidth}{0.5pt}}

% Dates section
%\section*{Dates:}

\end{document}
